%% kapa: exam=1 round=32 year=2021 period=1 form=A subject=civil_law subject_ko=민법 seq=9 number=9
%% source: Q-Net 감정평가사 기출문제 (cst00309) artlSeq=5208397 file=extracted/1008/5208397/A형/2021년도 제32회 감정평가사 1차 1교시 A형.pdf
%% conversion: kordoc 4.24.51 -> tools/gen1.py (auto); review=proofread-sonnet; flags=none
\begin{question}{9}
반사회적 법률행위에 관한 설명으로 옳지 않은 것은? (다툼이 있으면 판례에 따름)
\begin{choices}
  \item 어느 법률행위가 사회질서에 반하는지 여부는 특별한 사정이 없는 한 법률행위 당시를 기준으로 판단해야 한다.
  \item 강제집행을 면할 목적으로 부동산에 허위의 근저당권을 설정하는 행위는 특별한 사정이 없는 한 반사회적 법률행위라고 볼 수 없다.
  \item 대리인이 매도인의 배임행위에 적극 가담하여 이루어진 부동산의 이중매매의 경우, 본인인 매수인이 그러한 사정을 몰랐다면 반사회적 법률행위가 되지 않는다.
  \item 법률행위의 성립과정에서 단지 강박이라는 불법적 방법이 사용된 것에 불과한 때에는 반사회적 법률행위로 볼 수 없다.
  \item 반사회적 법률행위임을 이유로 하는 무효는 선의의 제3자에게 대항할 수 있다.
\end{choices}
\end{question}
\answer{3}
